% B_workflow.tex
\chapter{Scripts and workflow}\label{sec:workflow}

Every number and every geometry in this report is produced by a script that can be run again, and every script checks its inputs against values computed independently before it writes anything. This appendix describes the chain by function.

\section{Tools and machines}\label{sec:wf-tools}

\begin{itemize}
\item \textbf{Geometry}: Python with CadQuery (OpenCASCADE), writing STEP solids for the design regions, fixed regions, keep-outs and component envelopes from one parameter set.
\item \textbf{Analysis and optimisation}: Abaqus/CAE 2022 on a laptop (6 cores, \SI{16}{GB}) to build the model and write the solver and optimiser input; Abaqus and Tosca Structure 2024 on a server with 64 cores and \SI{251}{GB} for the production runs, with Abaqus 2025 on a second server for the coarse runs.
\item \textbf{Post-processing}: Abaqus Python on the server for the result databases, which the 2022 interface cannot open; NumPy and SciPy, with gmsh and manifold3d for meshing and Boolean operations on meshes, for iso surfaces, restoration, wall thickness and section properties.
\item \textbf{Figures}: matplotlib for plots, from tabulated data in the repository; Blender (Cycles) for the renders, built from the printed meshes and from component models drawn to data-sheet dimensions.
\item \textbf{Report}: \LaTeX{} with a shared preamble for both reports, packed into self-contained archives for an online editor.
\end{itemize}

\section{The chain}\label{sec:wf-chain}

\begin{enumerate}
\item \textbf{Sizing.} A script evaluates the propulsion model from the bench table: the thrust fit, the four power models and their physical test, hover point, flight time, sensitivities, sag and top speed. The loads script turns its outputs and the masses into the load cases.
\item \textbf{Material.} A script computes the merit indices from the data sheets, applies the gates and builds the design values.
\item \textbf{Hand studies.} Beam scripts with an energy method estimate the arm frequency for different profiles and tapers, and a joint script sizes bolts, fin and prongs in bearing.
\item \textbf{Geometry generation.} One generator writes the tub, lid and arm solids, split into design and fixed regions, the keep-outs and the component envelopes, with an export record of every volume. Separate generators write the keels and the retainers from the same parameters.
\item \textbf{Geometry checks at export.} A gate script runs the file-identity, Boolean bookkeeping, face-quality, tangency, test-meshing and clearance checks; nothing is exported with a failure. A second script computes the values the analysis model must reproduce, so the model build can be checked against them.
\item \textbf{Model build.} The model is built in Abaqus/CAE from a step-by-step set-up sheet, partitioning before any surface is selected, and every set, surface, mass and tie is checked against the expected values before the input file is written.
\item \textbf{Optimisation.} The solver and optimiser input files are transferred to the server and run there in thread mode, with every cycle's frequencies, compliances and support reactions reported.
\item \textbf{Thresholding and verification.} A script rewrites the optimiser's last cycle as a plain model with every element above the iso level solid and every other element near air, removes the optimiser's response definitions, and runs the frequency and static steps. Island checks run per part.
\item \textbf{Print geometry.} Marching tetrahedra on the nodal densities give the iso surface, Taubin smoothing removes facet noise, and Boolean operations restore the exact fixed solids and trim to the envelope. Print-geometry checks confirm one piece per part, empty keep-outs, no overlap between parts, and the wall thickness by ray casting.
\item \textbf{Strength and capacity.} A script reads the stresses of the final thresholded model per part and load case, applies the acceptance rule, and lists clusters and peaks; a section script cuts the printed arm and fin and computes their section moduli.
\item \textbf{Figures and report.} The plot script reads tabulated results; the render script builds the scenes; a pack script assembles the report archives.
\end{enumerate}

\section{Practices}\label{sec:wf-practices}

Every script that consumes a file checks its identity first: solid count, names and volumes against the export record. Every analysis prints its checks against expected values and stops with a clear message on a mismatch. Result databases from the optimiser are flattened to one instance, so post-processing searches instance-level sets. Geometry and results are kept under version control with a written note for every export, so each figure and number in the report can be traced to the run that produced it.
